\begin{textsample}{POS Dim 1 – 2011\_c4 – Score -3.00 – 2011\_c4/enade\_2011\_ef\_5f734ad795.txt}  \label{ex:f1_pos_322}
CAPÍTULO 4 PERCEPÇÃO DA PROVA As análises feitas neste capítulo tratam das percepções dos concluintes da Área de Educação Física ( Licenciatura ) sobre a prova aplicada no ENADE/2011.
Estas percepções foram mensuradas por meio de nove questões que avaliaram desde o grau de dificuldade da prova até o tempo gasto para concluí -la.
As percepções sobre a prova foram relacionadas com o desempenho dos estudantes e com a Grande Região de funcionamento do curso.
O questionário de percepção da prova encontra -se ao final do Anexo IV, que traz a reprodução da prova.
O desempenho dos estudantes foi classificado em quatro quartos.
Para tanto, esse desempenho foi ordenado de forma ascendente.
O percentil 25, P25, também conhecido como primeiro quartil, é a nota de desempenho que deixa um quarto ( 25 \% ) dos valores observados abaixo e três quartos acima.
A Figura 1 apresenta uma ilustração deste conceito.
O quarto inferior de desempenho é composto pelas notas abaixo do primeiro quartil.
Já o percentil 75, P75, também conhecido como terceiro quartil, é o valor para o qual há três quartos ( 75 \% ) dos dados abaixo e um quarto acima dele.
O quarto superior de desempenho é composto pelas notas iguais ou acima do terceiro quartil.
O percentil 50, P50, também conhecido como mediana, é o valor que divide as notas em dois conjuntos de igual tamanho.
O segundo quarto inclui valores entre o primeiro quartil ( P25 ) e a mediana.
O terceiro quarto contém os valores entre a mediana ( P50 ) e o terceiro quartil ( P75 ).
Vale ressaltar que percentis, quartis e medianas são pontos que não obrigatoriamente pertencem ao conjunto original de dados, ao passo que os quartos são subconjuntos dos dados originais.
A seguir, serão apresentados gráficos com resultados selecionados, relativos às nove questões avaliadas por grupos de estudantes.
Os gráficos apresentam nas barras o percentual de alunos que assinalaram uma das opções ou a soma das porcentagens daqueles que assinalaram duas ( ou três ) delas.
Por exemplo, para as questões 1 e 2, os gráficos apresentam a porcentagem total de participantes que assinalaram as opções ( D ) difícil e ( E ) muito difícil.
Em cada barra foram assinalados também os extremos do intervalo de confiança de 95 \% como linhas verticais unidas por uma linha horizontal na forma da letra H maiúscula, semelhantemente aos gráficos do Capítulo 3.
As Tabelas no Anexo II apresentam os valores absolutos e a distribuição percentual das alternativas válidas das nove questões, segundo o mesmo recorte de desempenho dos alunos e Grande Região de funcionamento do curso. 4.1 GRAU DE DIFICULDADE DA PROVA 4.1.1 Componente de Formação Geral Ao avaliarem “ Qual o grau de dificuldade desta prova na parte de Formação Geral? ” ( Questão 1 ), 26,5 \% do grupo de inscritos e presentes optaram pelas alternativas difícil ou muito difícil.
Entretanto, para mais da metade dos estudantes ( 61,3 \% ), o Componente de Formação Geral da prova foi considerado com grau de dificuldade médio ( Gráfico 4.1, Gráfico 4.2 e, no Anexo II, a Tabela II.1 ).
O percentual de estudantes que consideraram a prova como difícil ou muito difícil foi maior na região Sul, onde a proporção foi de 30,5 \%, enquanto a de menor incidência foi a região Nordeste, com 16,9 \%.
No Gráfico 4.1 é possível observar que a proporção encontrada na região Nordeste é menor do que a de todas as demais regiões com diferença estatisticamente significativa.
Nas Grandes Regiões, a proporção de presentes à prova que consideraram o Componente de Formação Geral como sendo de grau de dificuldade médio esteve entre 59,6 \% na região Sul e 67,5 \% na região Nordeste.
O percentual de alunos que consideraram a prova difícil ou muito difícil foi igual nos dois quartos de menor desempenho ( 30,6 \% ).
Nos quartos de maior desempenho a proporção de alunos que consideraram a prova difícil ou muito difícil – 25,7 \% no 3º quarto e 19,7 \% no 4º quarto – foi menor do que nos demais e significativamente diferentes, numa tendência decrescente.
Para todos os quartos de desempenho a alternativa modal para esta pergunta foi médio, com respectivamente 54,4 \% e 64,9 \% dos respondentes nos quartos extremos ( primeiro e último ). 4.1.2 Componente de Conhecimento Específico Ao responderem à Questão 2 – “ Qual o grau de dificuldade desta prova na parte de Componente Específico? ” – 24,4 \% do grupo de estudantes classificaram -na como difícil ou muito difícil.
Além disso, o Componente de Conhecimento Específico da prova foi considerado com grau de dificuldade médio por 61,6 \% dos alunos ( Gráfico 4.3, Gráfico 4.4, e, no Anexo II, a Tabela II.2 ).
A análise das respostas dos estudantes quanto ao grau de dificuldade do Componente de Conhecimento Específico da prova, agregados por Grande Região, mostra que a diferença entre a maior e a menor proporção de alunos que a avaliaram como difícil ou muito difícil é estatisticamente significativa : 27,8 \% na região Sul e 17,3 \% na Nordeste.
O percentual de alunos que classificaram o grau de dificuldade como médio, no Componente de Conhecimento Específico, variou de 60,4 \% a 65,5 \%, para as regiões, Sul e Norte, respectivamente.
Considerando -se a avaliação da dificuldade das questões do Componente de Conhecimento Específico da prova, de acordo com o desempenho dos estudantes, não se observa diferença estatisticamente significativa entre os quartos de menor desempenho : 27,7 \% no primeiro quarto e 26,7 \% no segundo.
Nos quartos de desempenho superiores, a proporção dos que classificaram a parte específica como difícil ou muito difícil foi de 20,0 \% ( último quarto ) e 23,4 \% ( terceiro quarto ).
A alternativa modal para a Questão 2 foi o grau médio, com 55,9 \% do quarto inferior e 63,3 \% do superior optando por esta resposta. 4.2 EXTENSÃO DA PROVA EM RELAÇÃO AO TEMPO TOTAL Indagados quanto à extensão da prova, em relação ao tempo total oferecido para a sua resolução ( Questão 3 ), os estudantes apontaram, com maior incidência, a alternativa que considerava a extensão adequada, para todas as agregações consideradas ( Gráfico 4.5, Gráfico 4.6, e, no Anexo II, a Tabela II.3 ).
O percentual de alunos que responderam ser a extensão da prova adequada foi de 51,6 \%.
Já 43,9 \% dos inscritos presentes consideraram que o exame foi longo ou muito longo e menos do que 5 \% o avaliaram como curto ou muito curto.
Entre as Grandes Regiões a proporção daqueles que avaliaram a prova como longa ou muito longa em relação ao tempo total destinado à sua resolução variou de 34,6 \% na região Norte até 45,8 \% na região Sul.
A diferença da região Norte é estatisticamente significativa em relação às regiões Nordeste, Sudeste e Sul.
Considerando -se o desempenho dos alunos, o percentual modal também foi considerar a extensão da prova adequada, em todos os quartos.
Para os de desempenho inferior, 46,9 \% assim consideraram e para os do quarto superior foram 54,1 \%.
No Gráfico 4.6, pode -se constatar que para os quatro quartos de desempenho a proporção de estudantes que consideraram a prova longa ou muito longa, em relação ao tempo total destinado à sua resolução, foi de aproximadamente 42 \% nos quartos superiores, de 44,5 \% no segundo quarto e de 46,7 \% no primeiro, sem diferenças estatisticamente significativas.
Existe, aparentemente, uma tendência decrescente com os quartos de renda. 4.3 COMPREENSÃO DOS ENUNCIADOS DAS QUESTÕES 4.3.1 Componente de Formação Geral Com relação aos enunciados das questões do Componente de Formação Geral ( Questão 4 ), as opiniões foram positivas, já que 73,3 \% dos alunos avaliados consideraram os enunciados de todas ou da maioria das questões claros e objetivos ( Gráfico 4.7, Gráfico 4.8, e, no Anexo II, a Tabela II.4 ).
Na análise regional, a percentagem de estudantes que avaliaram que todos ou a maioria dos enunciados das questões do Componente de Formação Geral estavam claros e objetivos variou de 71,2 \% na região Sudeste a 77,9 \% na região Centro-Oeste, sendo esta diferença estatisticamente significativa.
A análise das percepções dos estudantes sobre a clareza e objetividade dos enunciados permite afirmar que todos ou a maioria dos enunciados de questões relativas ao Componente de Formação Geral foram considerados claros e objetivos para a maior parte dos respondentes ( maior do que 71 \% em todas as regiões e maior do que 61 \% para todos os quartos de desempenho ).
Segundo o desempenho, observa -se que a proporção dos que emitiram esta opinião cresce conforme o desempenho aumenta, com diferenças estatisticamente significativas.
No quarto superior, a clareza e objetividade de todos ou da maioria dos enunciados das questões foi percebida por 81,8 \% dos alunos e no quarto de desempenho inferior tal avaliação foi emitida por 61,7 \% deles. 4.3.2 Componente de Conhecimento Específico Com relação aos enunciados das questões do Componente de Conhecimento Específico da prova, para 77,3 \% dos estudantes avaliados da Área de Educação Física ( Licenciatura ) a clareza e a objetividade ( Questão 5 ) estavam presentes em todas ou na maioria das questões ( Gráfico 4.9, Gráfico 4.10, e no Anexo II, a Tabela II.5 ).
A maioria dos estudantes de todas as Grandes Regiões brasileiras considerou claros e objetivos todos ou a maioria dos enunciados das questões do Componente de Conhecimento Específico da prova, percentual sempre maior do que 75 \%.
A região Sudeste foi onde ocorreu a menor proporção ( 75,4 \% ), com diferença estatisticamente significativa em relação a todas as outras regiões.
A proporção de estudantes que consideraram os enunciados das questões claros e objetivos apresenta uma tendência crescente em relação ao aumento de desempenho : mais elevada no quarto superior ( 85,4 \% ) se comparada ao quarto inferior de desempenho ( 65,1 \% ).
E esta diferença é estatisticamente significativa. 4.4 SUFICIÊNCIA DAS INFORMAÇÕES/INSTRUÇÕES FORNECIDAS Ao avaliarem as informações/instruções fornecidas para a resolução das questões ( Questão 6 ), 87,0 \% dos respondentes da Área de Educação Física ( Licenciatura ) de todo o Brasil afirmaram que estas eram até excessivas ou suficientes em todas ou na maioria das questões ( Gráfico 4.11, Gráfico 4.12, e, no Anexo II, a Tabela II.6 ).
Quanto à distribuição de respondentes pelas Grandes Regiões observa -se que a proporção de estudantes que consideraram as informações/instruções fornecidas até excessivas ou suficientes em todas ou na maioria das questões foi sempre superior a 85,0 \%, chegando a 88,3 \% na região Sul e a 88,5 \% na região Centro-Oeste.
As diferenças entre as regiões não chegaram a ser estatisticamente significativas.
Levando -se em conta o desempenho dos participantes, notam -se diferenças estatisticamente significativas entre as opiniões de estudantes a cada par de quartos de desempenho, como mostra o Gráfico 4.12.
O percentual de participantes que avaliaram as informações/instruções como até excessivas ou suficientes em todas ou na maioria das questões foi mais elevado no quarto superior ( 92,7 \% ), percentual superior à média nacional ( 87,0 \% ).
Já no quarto inferior, a suficiência das informações/instruções declarada como até excessiva em todas ou na maioria das questões foi percebida por 79,5 \% dos respondentes. 4.5 DIFICULDADE ENCONTRADA AO RESPONDER À PROVA Perguntados sobre as dificuldades com as quais se depararam ao responder à prova ( Questão 7 ), 8,9 \% dos estudantes apontaram o desconhecimento do conteúdo.
Para 41,2 \%, a forma diferente de abordagem do conteúdo foi indicada como dificuldade.
Já a falta de motivação para fazer a prova foi a dificuldade apontada por 23,3 \% dos respondentes.
Considerando -se todo o Brasil, 22,8 \% dos respondentes afirmaram que não tiveram qualquer tipo de dificuldade para responder à prova ( Tabela II.7 no Anexo II ).
Os Gráficos 4.13 e 4.14 apresentam os percentuais de estudantes que apontaram o desconhecimento do conteúdo como dificuldade percebida ao responder à prova.
Na análise por Grandes Regiões, o percentual de inscritos e presentes que apontaram o desconhecimento do conteúdo como dificuldade ao responder à prova variou de 7,8 \% na região Nordeste a 10,1 \% na Centro-Oeste, sendo que esta diferença não é significativa estatisticamente.
A forma diferente de abordagem do conteúdo foi a escolha modal dos estudantes, com percentuais que variaram de 33,1 \% ( região Nordeste ) a 44,6 \% ( Norte ).
O percentual de alunos que citou a falta de motivação como dificuldade variou de 19,7 \% ( região Sul ) a 29,8 \% ( região Nordeste ).
Os que declararam não ter qualquer dificuldade para responder à prova variaram de 22,1 \% na região Centro-Oeste a 24,8 \% na Nordeste.
Com relação aos quartos de desempenho, o desconhecimento do conteúdo, foi a opção escolhida por 8,0 \% dos estudantes do quarto superior e 10,8 \% do quarto inferior.
A alternativa modal para os alunos, quando agregados pelos quartos de desempenho, foi que a dificuldade encontrada foi causada pela forma diferente de abordagem do conteúdo : 37,0 \% no quarto inferior e 40,2 \% do quarto superior assim o responderam. 4.6 CONTEÚDOS DAS QUESTÕES OBJETIVAS DA PROVA Ao analisarem os conteúdos das questões objetivas da prova ( Questão 8 ), um percentual muito pequeno dos estudantes avaliados, apenas 4,8 \%, afirmou que não estudou ainda a maioria desses conteúdos ( Gráficos 4.15, Gráfico 4.16, e a Tabelas II.8 no Anexo II ).
A maioria ( 74,0 \% ) afirmou ter estudado e aprendido muitos ou todos os conteúdos.
Na análise por Grande Região, a proporção de respondentes que escolheram a opção não estudou ainda a maioria desses conteúdos, foi pequena.
Observa -se que nas regiões Nordeste ( 4,9 \% ), Sudeste ( 5,2 \% ) e Centro-Oeste ( 5,6 \% ), apesar de pequenas, as proporções foram maiores do que a média nacional ( 4,8 \% ).
Não se observa diferença estatisticamente significativa entre as regiões.
Em todas as regiões, a maioria dos presentes afirmou ter estudado e aprendido muitos ou todos os conteúdos, com proporções variando entre 70,5 \% na região Centro-Oeste e 76,1 \% na Norte.
Considerando -se separadamente as opiniões de estudantes dos quatro quartos de desempenho, observa -se que, no quarto inferior, 10,9 \% ofereceram como resposta que não estudou ainda a maioria desses conteúdos, enquanto no quarto superior apenas 1,3 \% deram a mesma resposta, uma proporção decrescente com os quartos de desempenho.
A diferença entre os alunos que optaram por este motivo de dificuldade é estatisticamente significativa entre os quartos.
Tendo em conta o quarto superior, 87,9 \% dos alunos afirmaram ter estudado e aprendido muitos ou todos os conteúdos. 4.7 TEMPO GASTO PARA CONCLUIR A PROVA Ao responderem sobre o tempo de conclusão da prova ( Questão 9 ), mais da metade dos estudantes ( 64,4 \% ) afirmou ter gasto entre duas e quatro horas ( Gráfico 4.17, Gráfico 4.18 e, no Anexo II, a Tabela II.9 ).
Considerando -se as cinco Grandes Regiões brasileiras, o percentual dos que utilizaram entre duas e quatro horas para finalizar a prova, apenas na região Nordeste ( 60,5 \% ) foi inferior ao percentual nacional ( 64,4 \% ).
Nas demais regiões o percentual de alunos que dispensaram entre duas e quatro horas para concluir a prova variou de 64,7 \% ( região Sudeste ) a 67,5 \% ( região Norte ), como mostra o Gráfico 4.17.
Uma vez sendo analisadas as alternativas escolhidas pelos estudantes agrupados por quartos de desempenho, observa -se que uma maior proporção de participantes no quarto superior declarou ter gasto entre duas e quatro horas para concluir a prova quando comparados com os do quarto inferior, respectivamente 76,7 \% e 43,7 \%.
As diferenças entre os quartos são estatisticamente significativas e evidenciam uma tendência crescente.

% matched lemmas: 
\end{textsample}
